The original Transformer relies on two optimization details that are each described in a single sentence: a
learning-rate warmup and the scaling of attention logits by $1/\sqrt{d_k}$. Both are copied into most
implementations, yet it is rarely tested which of them a given model needs. We ablate both factors in a controlled
$2\times 2$ design on a small post-LayerNorm encoder--decoder Transformer trained to reverse digit sequences, with
\valSeeds{} seeds per configuration. Without warmup, every run fails to learn: exact-match accuracy is
\valAccNoWarmup\% and the training loss stays near \valPlateauLoss. A control that removes only the ramp, starting at the
warmup's peak learning rate, solves the task in \valRampSolved{} of \valRampRuns{} seeds, so the failure is not explained by
the large initial learning rate alone. Removing the attention scaling has no measurable
effect on the outcome: accuracy is \valAccNoScaling\% without it and \valAccBaseline\% with it. At this scale ($d_k=\valDk$), warmup
matters and attention scaling does not; whether the scaling matters for larger key dimensions remains untested.
